\section*{Bab \ref{ch:b3:advanced-geochronology} --- Geokronologi Lanjut}

\begin{solution}{exo:b3:advanced-geochronology:1}
$t_{206} = \ln(1.04500)/\lambda_{238} = \qty{283.8}{\Ma}$; $t_{207} =
\ln(1.3300)/\lambda_{235} = \qty{289.6}{\Ma}$. Selisih kedua umur sebesar
\qty{6}{\Ma}: zirkon itu diskordan. $t_{207}$-nya lebih tua, sehingga titik itu jatuh
di bawah \omterm{def:b3:advanced-geochronology:concordia}{konkordia}, di sisi hilangnya timbal.
\end{solution}

\begin{solution}{exo:b3:advanced-geochronology:2}
$J = (e^{5.543\times10^{-10}\times2.8201\times10^7} - 1)/1.3129 =
0.015754/1.3129 = 0.01200$. Maka $t = \ln(1 + 0.01200\times10.000)/\lambda =
\ln 1.12/(5.543\times10^{-10}) = \qty{204.5}{\Ma}$.
\end{solution}

\begin{solution}{exo:b3:advanced-geochronology:3}
$T_{\mathrm{CHUR}} = \ln\bigl(1 + (0.51180 - 0.512638)/(0.1200 -
0.1966)\bigr)/\lambda = \ln(1.010940)/\lambda = \qty{1.67}{\Ga}$;
$T_{\mathrm{DM}} = \ln\bigl(1 + (0.51180 - 0.51315)/(0.1200 - 0.2137)\bigr)/\lambda
= \ln(1.014408)/\lambda = \qty{2.19}{\Ga}$.
\end{solution}

\begin{solution}{exo:b3:advanced-geochronology:4}
Tanpa peluruhan, $t = N/P_0 = \num{2.0e5}/4.96 = \qty{40.3}{ka}$. Dengan peluruhan,
$t = -\ln(1 - N\lambda/P_0)/\lambda = -\ln(1 - 0.02014)/\lambda =
\qty{40.7}{ka}$, \qty{1}{\percent} lebih tua.
\end{solution}

\begin{solution}{exo:b3:advanced-geochronology:5}
Untuk $t$ kecil, $x \approx \lambda_{235}t$ dan $y \approx \lambda_{238}t$, sehingga
$y \approx (\lambda_{238}/\lambda_{235})x$, kemiringannya $0.1575$. Jika timbal
hilang baru-baru ini, $t_2 \approx 0$ dan \omterm{def:b3:advanced-geochronology:concordia}{titik potong bawah} adalah $C(0)$, yaitu
titik asal. Sepanjang garis yang melalui titik asal, $y/x$ konstan, begitu pula
$^{207}\mathrm{Pb}^*/^{206}\mathrm{Pb}^* = x/(137.8\,y)$: setiap kristal memberikan
umur timbal--timbal pada \omterm{def:b3:advanced-geochronology:concordia}{titik potong atas}, yaitu umur pembentukan.
\end{solution}

\begin{solution}{exo:b3:advanced-geochronology:6}
Bobot $1.5625$, $1$, $1.2346$, $2.0408$, $W = 5.838$; $\bar x =
\qty{296.09}{\Ma}$, $\sigma = 1/\sqrt W = \qty{0.41}{\Ma}$. Simpangan kuadrat terbobot
berjumlah $2.32$; dengan $f = 3$, $\mathrm{MSWD} = 0.77$, berada di bawah
$1 + 2\sqrt{2/3} = 2.63$: umur-umurnya konsisten.
\end{solution}

\begin{solution}{exo:b3:advanced-geochronology:7}
$\lambda\zeta g\rho_d\rho_s/\rho_i = 1.55125\times10^{-10}\times350\times0.5
\times\num{1.10e6}\times0.5 = 0.014931$; $t = \ln(1.014931)/\lambda =
\qty{95.5}{\Ma}$. Jika batuan mendingin cepat, umur ini mencatat saat batuan melewati
zona penganilan parsial apatit, sekitar \qty{60}{} hingga \qty{120}{\celsius}; jika
mendingin lambat, umurnya campuran, dan panjang jejaklah yang harus menjelaskan.
\end{solution}

\begin{solution}{exo:b3:advanced-geochronology:8}
Per gram: uranium $20\times10^{-6}/238.05\times N_A = \num{5.06e16}$ atom, dengan
\num{5.02e16} uranium-238 dan \num{3.64e14} uranium-235; torium
\num{1.038e17}, $\lambda_{232} = \ln 2/\num{1.40e10} = \qty{4.95e-11}{yr^{-1}}$;
helium $10^{-9}N_A = \num{6.02e14}$. \omterm{def:b3:advanced-geochronology:cosmogenic}{Laju produksi} adalah
$8\lambda_{238}{}^{238}\mathrm U + 7\lambda_{235}{}^{235}\mathrm U +
6\lambda_{232}\mathrm{Th} = (6.23 + 0.25 + 3.08)\times10^7 =
\num{9.57e7}\text{ atom}\,\unit{g^{-1}.yr^{-1}}$, dan $t = \num{6.02e14}/\num{9.57e7} =
\qty{6.29}{Ma}$. Terkoreksi: $6.29/0.75 = \qty{8.39}{\Ma}$.
\end{solution}

\begin{solution}{exo:b3:advanced-geochronology:9}
$P_0/\lambda = 4.96/\num{4.994e-7} = \num{9.93e6}\text{ atom}\,\unit{g^{-1}}$, tercapai hingga
\qty{90}{\percent} setelah $\ln 10/\lambda = \qty{4.6}{Ma}$. Pada
\qty{10}{m.Ma^{-1}} $= \qty{1e-3}{cm.yr^{-1}}$: $\rho\varepsilon/\Lambda =
2.7\times10^{-3}/160 = \qty{1.69e-5}{yr^{-1}}$, $\mu = \qty{1.74e-5}{yr^{-1}}$
dan $N = 4.96/\num{1.74e-5} = \num{2.85e5}\text{ atom}\,\unit{g^{-1}}$, 35 kali lebih sedikit.
\end{solution}

\begin{solution}{exo:b3:advanced-geochronology:10}
Dengan $W = w_1 + w_2$, $x_1 - \bar x = w_2(x_1 - x_2)/W$ dan $x_2 - \bar x =
-w_1(x_1 - x_2)/W$, sehingga $\sum w_i(x_i - \bar x)^2 = (x_1 - x_2)^2(w_1w_2^2 +
w_2w_1^2)/W^2 = (x_1 - x_2)^2w_1w_2/W = (x_1 - x_2)^2/(\sigma_1^2 +
\sigma_2^2)$; dengan $f = 1$ ini adalah MSWD. Jika keduanya mengukur umur yang sama,
$x_1 - x_2$ memiliki rata-rata 0 dan varians $\sigma_1^2 + \sigma_2^2$, sehingga
harapan kuadratnya dibagi varians itu bernilai 1.
\end{solution}

\begin{solution}{exo:b3:advanced-geochronology:11}
$C(2700) = (13.283, 0.52019)$, $C(500) = (0.63628, 0.08065)$. $f = (0.52019 -
0.300)/(0.52019 - 0.08065) = 0.501$, setengah timbalnya. $x = 0.499\times13.283 +
0.501\times0.63628 = 6.948$. $t_{206} = \qty{1691}{\Ma}$, $t_{207} =
\qty{2105}{\Ma}$; rasio $x/(137.8\,y) = 0.1681$ secara iteratif memberikan
umur timbal--timbal sebesar \qty{2537}{\Ma}. Tidak satu pun dari ketiganya bermakna:
campuran dari kedua peristiwa, dan hanya titik potong yang menentukan umurnya. Umur timbal--timbal
merupakan nilai minimum bagi kristalisasi.
\end{solution}

\begin{solution}{exo:b3:advanced-geochronology:12}
Tanpa erosi: $t = -\ln(1 - \num{1.0e6}\times\num{4.994e-7}/4.96)/\lambda =
\qty{212}{ka}$. Keadaan tunak: $\rho\varepsilon/\Lambda = P_0/N - \lambda =
\qty{4.46e-6}{yr^{-1}}$, $\varepsilon = \qty{2.64e-4}{cm.yr^{-1}} =
\qty{2.6}{m.Ma^{-1}}$. Untuk sembarang $\varepsilon$ di bawah ini, $N =
P_0(1 - e^{-\mu t})/\mu$ menentukan satu $t$, mulai dari \qty{212}{ka} naik hingga
tak hingga: satu pengukuran memiliki dua besaran tak diketahui. Nuklida kedua dengan
waktu paruh berbeda, misalnya aluminium-26 dalam kuarsa yang sama, atau konsentrasi pada
kedalaman dalam inti bor, memisahkannya.
\end{solution}

\begin{solution}{pb:b3:advanced-geochronology:1}
\textbf{1.} $t_{206}$ / $t_{207}$, dalam \unit{\Ma}: Z1 301.9 / 302.2, Z2 284.1 /
286.5, Z3 265.7 / 268.5, Z4 246.1 / 250.1, Z5 225.1 / 229.4. Galat sebesar
\qty{0.2}{\percent} pada $y$ kira-kira \qty{0.6}{\Ma}: Z1 konkordan dalam
galat; yang lain diskordan, dengan derajat yang makin besar dari Z2 ke Z5.

\textbf{2.} Zirkon-zirkon itu kehilangan fraksi timbal yang makin besar dalam satu
peristiwa berikutnya: titiknya makin jauh ke bawah tali busur. Tali busur berada di bawah
\omterm{def:b3:advanced-geochronology:concordia}{konkordia} yang cekung, tempat $x$ lebih besar untuk $y$ yang sama: $t_{207} > t_{206}$.

\textbf{3.} Z1: $X = 1/0.04794 = 20.859$, $Y = 0.3467/(137.8\times0.04794) =
0.05248$. Z5: $X = 28.145$, $Y = 0.05178$.

\textbf{4.} Bobot $1.000$, $1.235$, $1.563$, $1.563$, $2.041$ (dalam
$10^8$); $\bar x = 0.29402$, $\bar y = 0.040892$, $S_{xx} = \num{7.845e5}$.

\textbf{5.} $b = 0.13260$, $a = \bar y - b\bar x = 0.001905$.

\textbf{6.} Pada \qty{310}{\Ma}, $\mathrm dy/\mathrm dx = 0.1575\,
e^{(\lambda_{238} - \lambda_{235})t} = 0.1218$. \omterm{def:b3:advanced-geochronology:concordia}{Diskordia} hanya
\qty{9}{\percent} lebih curam daripada \omterm{def:b3:advanced-geochronology:concordia}{konkordia}: ia bertemu dengannya pada sudut kecil,
dan \omterm{def:b3:advanced-geochronology:concordia}{titik potong atas} akan kurang presisi.

\textbf{7.} Newton dari \qty{310}{\Ma}: $t = \qty{311.2}{\Ma}$.

\textbf{8.} Dari \qty{100}{\Ma}: $t = \qty{99.7}{\Ma}$.

\textbf{9.} $(y_i - a - bx_i)/\sigma_i = 0.64$, $-0.91$, $0.35$, $-0.24$,
$0.17$; $\mathrm{MSWD} = 1.435/3 = 0.48$, di bawah batas $2.63$ untuk $f = 3$:
zirkon-zirkon berada pada satu garis.

\textbf{10.} $s_b = 1/\sqrt{\num{7.845e5}} = 0.00113$.

\textbf{11.} Dengan $b - 2s_b = 0.13034$ \omterm{def:b3:advanced-geochronology:concordia}{titik potong atas} adalah
\qty{322.9}{\Ma}; dengan $b + 2s_b = 0.13485$, \qty{302.6}{\Ma}. Umurnya
$311 \pm \qty{10}{\Ma}$ pada dua simpangan baku.

\textbf{12.} $f = (x(311.2) - x_i)/(x(311.2) - x(99.7))$: Z1 \qty{4.7}{\percent},
Z5 \qty{41}{\percent}.

\textbf{13.} $J = 0.015754/1.3129 = 0.0119999$, $0.01200$.

\textbf{14.} Langkah 1: \qty{150.0}{\Ma}; langkah 4: \qty{296.6}{\Ma}; langkah 10:
\qty{295.0}{\Ma}.

\textbf{15.} Langkah 4 sampai 10: tujuh langkah berurutan, \qty{82}{\percent} dari
argon-39, tanpa tren. Langkah 1 sampai 3 naik seperti tangga.

\textbf{16.} Galat sama: rata-rata $2072.9/7 = \qty{296.1}{\Ma}$, $\sigma =
1.1/\sqrt7 = \qty{0.42}{\Ma}$; jumlah kuadrat simpangannya $3.83$,
$\mathrm{MSWD} = 3.83/(1.21\times6) = 0.53$, di bawah $2.15$.

\textbf{17.} $\bar R = \sum f_iR_i = 14.270$, $t = \qty{285.2}{\Ma}$: argon
yang hilang dari tepinya, dilepaskan pada langkah-langkah awal, menurunkan rata-rata, sebagaimana
ia akan menurunkan umur kalium–argon.

\textbf{18.} $(700 - 375)/(311.2 - 296.1) = \qty{21.5}{\celsius.Ma^{-1}}$
(dengan rentang yang lebar, karena umur zirkon diketahui hingga $\pm\qty{10}{\Ma}$).

\textbf{19.} Sekitar \qty{100}{\Ma} granit itu terkubur di bawah batuan
penutup Mesozoikum, dan pemanasan atau sirkulasi fluida membuat zirkon-zirkon kehilangan timbal.
Biotit kehilangan argon dari tepinya, yang tercatat oleh bentuk tangga; langkah
terendahnya, \qty{150}{\Ma}, adalah campuran argon tua dan argon yang direset. \omterm{def:b3:advanced-geochronology:arar}{Umur plato}
tetap utuh: biotit tidak cukup dipanaskan untuk kehilangan seluruh argonnya.

\textbf{20.} Kemiringan kuadrat terkecil elevasi terhadap umur adalah
\qty{20.1}{m.Ma^{-1}}.

\textbf{21.} $(65 - 10)/30 = \qty{1.83}{km}$; sejak \qty{43.1}{\Ma}, rata-rata
$1830/43.1 = \qty{42}{m.Ma^{-1}}$.

\textbf{22.} $\rho\varepsilon/\Lambda = P_0/N - \lambda = \num{3.415e-5} -
\num{0.050e-5} = \qty{3.365e-5}{yr^{-1}}$; $\varepsilon = \num{3.365e-5}\times
160/2.7 = \qty{1.99e-3}{cm.yr^{-1}} = \qty{19.9}{m.Ma^{-1}}$.

\textbf{23.} $P_0\Lambda/(\rho N) = \qty{20.2}{m.Ma^{-1}}$, \qty{1.5}{\percent}
lebih besar: peluruhan hampir tidak berpengaruh pada laju ini.

\textbf{24.} $\Lambda/(\rho\varepsilon) = 160/(2.7\times\num{1.99e-3}) =
\qty{30}{ka}$. Granit naik sekitar \qty{20}{m.Ma^{-1}} antara 52 dan
\qty{43}{\Ma}, rata-rata \qty{42}{m.Ma^{-1}} sejak itu, dan kini terkikis dengan
\qty{20}{m.Ma^{-1}} sekarang: penyingkapan batuan penutup lebih cepat pada suatu saat di
Kenozoikum.

\textbf{25.} Granit Varden mengkristal
\textbf{\qty[text-series-to-math=true, separate-uncertainty=true]{311 \pm 10}{\Ma}} yang lalu (\omterm{def:b3:advanced-geochronology:concordia}{titik potong atas},
pada dua simpangan baku).
\end{solution}
